\chapter{Implementation}
\label{chap:implementation}

\section{Software Environment and Toolchain}
Development of the benchmarking suite required a balance between 
architectural modularity and strict low-level performance control. 
We selected C++20 for the engine specifically to use \texttt{std::span} 
and other modern primitives that provide memory safety without the runtime 
cost of managed pointers. 
To manage the five discrete order book implementations, we used the 
\texttt{CMake} build system, which allowed us to maintain a modular 
architecture while ensuring identical compiler flags across all variants.

Key components of the toolchain include:
\begin{itemize}
    \item \textbf{C++20 Primitive Types}: The system uses \texttt{std::span} 
    for non-owning views of message buffers. This choice was made to 
    eliminate unnecessary copies during FIX parsing while keeping 
    memory layouts consistent and predictable.
    \item \textbf{POSIX Socket API}: We opted for the raw \texttt{sys/socket.h} 
    API over higher-level networking libraries. By bypassing those 
    abstractions, we ensure the gateway's latency is purely a function 
    of the OS kernel stack traversal.
    \item \textbf{Functional Parity Testing}: To ensure that every order 
    book variant produces identical matching results, we built a 
    unit testing suite around the \texttt{GoogleTest} framework. This 
    validates that implementations remain functionally consistent, even 
    when their internal memory layouts differ.
\end{itemize}

\section{Order Book Implementations}
The core of this research is the implementation of five distinct order book 
strategies, each representing a different point on the efficiency-flexibility 
spectrum. This section details the technical implementation of the 
baseline, the static Array, the Hybrid "Hot-Cold" model, and the Memory Pool.

\subsection{Standard Associative Map (Baseline)}
The baseline implementation (\texttt{MapOrderBook}) uses \texttt{std::map} 
to link price levels with order queues. We use this to measure the 
latency cost of node-based allocations and pointer-chasing, as shown 
in the results in Table~\ref{tab:direct_mean_latency}.

\subsection{Direct-Indexed Array (O(1) Access)}
The \texttt{ArrayOrderBook} achieves O(1) price retrieval by mapping price 
levels directly to a contiguous array. While this approach is the fastest 
possible method for price lookups, it introduces a practical constraint: 
the book must know its price range at startup. We chose this trade-off 
to establish a theoretical performance floor for the study, fulfilling 
the requirements for low-latency execution (\FR{2}). Mapping prices 
to indices, as shown in Listing 4.1, is a simple arithmetic operation that 
remains stable regardless of the book's depth.

\begin{lstlisting}[language=C++, caption={O(1) Price-to-Index Mapping}]
// src/orderbooks/array_order_book.cpp

Index ArrayOrderBook::priceToIndex(Price price) const {
    // Normalise price to a zero-based index for the array
    return (price - minPrice_) / tickSize_;
}

void ArrayOrderBook::addOrder(const Order &order) {
    if (!isValidPrice(order.price)) return;

    // Direct access to the price level's memory location
    Index index = priceToIndex(order.price);
    auto &levels = (order.side == Side::Buy) ? bidLevels_ : askLevels_;
    
    levels[index].push_back(order);
    updateCaches(order); // O(1) cache update
}
\end{lstlisting}

This design achieves \FR{2} by providing deterministic latency, 
independent of book depth. However, it requires prior knowledge of 
the price range, a constraint common in fixed-range instruments.

\subsection{Hybrid "Hot-Cold" Vector-Map}
To balance the speed of contiguous memory with the flexibility of 
associative storage, the \texttt{HybridOrderBook} implements a tiered 
memory model. It maintains the "Hot" price levels (those closest to the 
spread) in a sorted \texttt{std::vector} for cache-friendly matching, while 
"Cold" levels (out-of-the-money) are stored in a \texttt{std::map}.

\begin{lstlisting}[language=C++, caption={Hot/Cold Proximity Check}]
// src/orderbooks/hybrid_order_book.cpp

void HybridOrderBook::addOrder(const Order &order) {
    // Check if the price level is within the performance "Hot Path"
    if (isCloseToSpread(order.price, order.side == Side::Buy)) {
        if (hotBids_.size() >= maxHotLevels_) {
            demoteFromHot(true); // Evict furthest level to Map
        }
        addToHot(order, true); // Insert into contiguous Vector
    } else {
        addToCold(order, true); // Store in associative Map
    }
}
\end{lstlisting}

This model balances the flexibility of a traditional order book (\FR{1}) 
with higher performance for price levels immediately surrounding the 
spread. The impact of the hot-path size on overall matching throughput 
is analyzed in Section~\ref{subsec:pool_hybrid_mech}.

\subsection{Pool-Allocated Order Book (Slab Allocation)}
Runtime heap allocations are a primary source of latency jitter in 
dynamic systems. The \texttt{PoolOrderBook} mitigates this by using 
a pre-allocated slab of \texttt{OrderNode} structures. By managing an 
internal free list, the engine can add or cancel orders without requesting 
memory from the OS kernel, ensuring the matching path remains 
deterministic.

\begin{lstlisting}[language=C++, caption={Slab-Allocation and Free List Management}]
// src/orderbooks/pool_order_book.cpp

Index PoolOrderBook::allocateSlot() {
    // O(1) Pop from the pre-allocated free list
    Index index = freeHead_;
    freeHead_ = orders_[index].nextFree;
    
    // Initialise node state without new/malloc
    orders_[index].next = NULL_IDX;
    orders_[index].prev = NULL_IDX;
    return index;
}

void PoolOrderBook::freeSlot(Index idx) {
    // Return node to pool for immediate recycling
    orders_[idx].nextFree = freeHead_;
    freeHead_ = idx;
}
\end{lstlisting}

By reusing memory slots, we keep the matching logic deterministic. This 
mitigates allocation jitter and helps verify the reduced tail latency 
analyzed in Section~\ref{subsec:pool_hybrid_mech}.

\section{The TCP Order Gateway}
We need to test the system under network conditions for H3, so we implemented 
a high-performance TCP gateway. This component connects external FIX 
clients to the internal core through a lock-free interface.

\subsection{Protocol Specification and Parsing}
Our FIX implementation focus was on parsing speed rather than full 
protocol coverage. The gateway uses a non-allocating incremental 
parser that processes raw buffers as they arrive. As Listing 4.3 
demonstrates, we used a sliding window technique to handle fragmented 
TCP packets, satisfying the realism requirement of \NFR{4} while keeping 
the matching core separate from network ingestion.

\begin{lstlisting}[language=C++, caption={Incremental FIX Parsing Loop}]
// src/network/tcp_order_gateway.cpp

void TCPOrderGateway::clientHandler(int clientSocket) {
    std::vector<char> buffer(4096);
    size_t offset = 0; // Tracks partial message data

    while (running_) {
        // Read data into the buffer starting after the leftover data
        ssize_t bytes = recv(clientSocket, buffer.data() + offset, 
                             buffer.size() - offset, 0);
        
        size_t totalBytes = offset + bytes;
        size_t processed = 0;

        while (processed < totalBytes) {
            size_t consumed = 0;
            std::span<const char> data(buffer.data() + processed, 
                                       totalBytes - processed);

            auto order = FIXParser::parse(data, consumed);
            if (consumed == 0) break; // Incomplete message, wait for more

            if (order) {
                order->receiveTimestamp = getCurrentTimeNs();
                while (!orderQueue_.push(*order)) std::this_thread::yield();
            }
            processed += consumed;
        }

        // Shift remaining data to the front for the next recv()
        if (processed < totalBytes) {
            std::memmove(buffer.data(), buffer.data() + processed, 
                         totalBytes - processed);
            offset = totalBytes - processed;
        } else {
            offset = 0;
        }
    }
}
\end{lstlisting}

\subsection{The Loopback Simulation Model}
We need to measure kernel networking overhead without the jitter of hardware 
NICs, so we run the gateway on the loopback interface (\texttt{127.0.0.1}). 
This allows us to isolate the "Wire-to-Engine" latency analyzed in 
Section~\ref{subsec:scenario_latency}.

\section{High-Performance Concurrency}
The system employs a multi-threaded architecture to separate high-latency 
I/O from the latency-sensitive matching core.

\subsection{Wait-Free SPSC Queue}
For Direct and single-client Gateway modes, we use a custom 
Single-Producer Single-Consumer (SPSC) queue. This implementation 
relies on a contiguous ring-buffer and atomic head/tail pointers 
with \texttt{std::memory\_order\_acquire/release} semantics, ensuring 
that the matching engine does not block in this design due to lock contention.

In multi-producer scenarios (\FR{4}), the system swaps the SPSC queue 
for the \texttt{ConcurrentQueue} library~\citep{moodycamel_concurrentqueue}. 
We opted for this established lock-free implementation to ensure our results 
reflect the underlying data structure performance rather than my own 
queue implementation logic. We evaluate this MPMC transport under 
contention in Section~\ref{sec:mpsc_contention}.

\section{Low-Level Performance Utilities}
To ensure nanosecond-level accuracy and deterministic execution across 
benchmarks, the system includes a set of low-level utilities for 
precise timing and thread management.

\subsection{Nanosecond Precision Timing}
To achieve the nanosecond granularity required for order book benchmarking, 
the system employs the \texttt{std::chrono} high-resolution clock. While 
cycle-level precision is typically associated with the CPU's Time Stamp 
Counter (\texttt{RDTSC}), our implementation uses the standard C++ library 
to ensure portability across environments (specifically WSL2) while 
maintaining the precision necessary for the evaluation in 
Chapter~\ref{chap:testing_evaluation}.

\begin{lstlisting}[language=C++, caption={Nanosecond Precision Timing}]
// src/utils/rdtsc.hpp

inline uint64_t getCurrentTimeNs() {
    // Cross-platform fallback. In production, this is replaced 
    // by __builtin_ia32_rdtsc() for cycle-level accuracy.
    return std::chrono::duration_cast<std::chrono::nanoseconds>(
               std::chrono::high_resolution_clock::now().time_since_epoch()
           ).count();
}
\end{lstlisting}

\subsection{Thread-to-Core Affinity (CPU Pinning)}
To eliminate the performance jitter caused by the Linux scheduler 
migrating threads across physical cores (the "Context Switching Jitter"), 
we implemented CPU pinning as shown in Listing 4.5.

\begin{lstlisting}[language=C++, caption={POSIX Thread Pinning Strategy}]
// src/utils/thread_pinning.hpp

inline bool pinToCore(int core_id) {
#if defined(__linux__)
    cpu_set_t cpuset;
    CPU_ZERO(&cpuset);       // Clear the CPU mask
    CPU_SET(core_id, &cpuset); // Set the bit for the target core

    // Request the OS scheduler to bind this thread to the selected core
    pthread_t thread = pthread_self();
    return pthread_setaffinity_np(thread, sizeof(cpu_set_t), &cpuset) == 0;
#endif
    return false; // Fallback for unsupported OS
}
\end{lstlisting}

By binding the matching engine thread to a specific physical core, we 
eliminate "Context Switching Jitter" where the OS scheduler migrates a 
thread between CPUs. This ensures that the engine's working set 
remains "hot" in the local L1/L2 caches, maintaining the deterministic 
execution time required to validate our performance claims. This pinning 
helps satisfy \NFR{4} by making sure our measurements are not skewed by 
operating system interference.

\section{Automation and Data Collection}
The suite uses a Python-based driver to sweep across the benchmark 
scenarios defined in Chapter~\ref{chap:system_design}. This driver coordinates 
the execution of each order book variant and manages the serialized CSV output 
required for the detailed evaluation in Chapter~\ref{chap:testing_evaluation}.
